\clearpage\begingroup
\HistoricalHead{168}{L}{PHYSIK.\quad BRIEFWECHSEL.}
\begin{center}
\textbf{Anhang zu [23.].}\par\medskip
\textbf{\textls[80]{[Preisaufgabe der Königlichen Gesellschaft der Wissenschaften.]}}\par\medskip
\rule{.76\linewidth}{.4pt}\par\smallskip
{\footnotesize Göttingische Gelehrte Anzeigen 1831, 196. Stück, S. 1947---1949.}\par\smallskip
\rule{.76\linewidth}{.4pt}
\end{center}
\noindent\hspace*{\parindent}Für den November d. J. [1831] war von der \textls[120]{mathematischen} Classe der Hauptpreis auf die Beantwortung der Aufgabe gesetzt:
\par\noindent . . . . . . . . . . . . . . . . . . . . .\par
\noindent\hspace*{\parindent}In der praktischen Astronomie mangelt es noch immer an einem Mittel zur sichern Bestimmung der Lichtstärke der Himmelskörper, und die früher zu diesem Zwecke in Vorschlag gebrachten Vorrichtungen haben sich in der Anwendung wenig brauchbar gezeigt.

Da es jedoch von vielfachem und grossem Nutzen seyn würde, die verschiedenen Abstufungen des Sternenlichts und die darin statt findenden Veränderungen mit Sicherheit und Leichtigkeit beurtheilen zu können:
\par\begingroup\leftskip=\parindent\parindent=0pt
\noindent so wünscht die Königliche Societät neue, durch vollständige Beschreibungen erläuterte Vorschläge zu solchen auf photometrischen Grundsätzen beruhenden Vorrichtungen zu erhalten, mittelst welcher die verschiedenen Grade des Lichts der Fixsterne mit Sicherheit, Gleichförmigkeit und Leichtigkeit beurtheilt und festgestellt werden können, und deren Leistungen aus einer ausführlichen Darlegung der Resultate, die aus ihrer Anwendung auf Sterne von den verschiedensten Grössen erhalten worden sind, sich erkennen und beurtheilen lassen.\par\endgroup
\noindent Es waren darauf zwey Concurrenzschriften mit den Denksprüchen\par
\noindent\hspace*{\parindent}\begin{tabular}{@{}l@{\quad}l@{}}
No. 1.&Errare humanum est;\\
No. 2.&Per aspera ad astra
\end{tabular}\par
\noindent bezeichnet eingegangen.
\clearpage\SourceLeaf{174}\HistoricalHead{169}{R}{PREISAUFGABE ÜBER PHOTOMETER.}
\noindent\hspace*{\parindent}Der Verfasser der ersten Abhandlung legt zuerst das Prinzip dar, nach welchem allein, seiner Meinung zufolge, dass Verhältniss der Lichtstärke der Himmelskörper bestimmt werden könne, gibt hierauf eine ausführliche, durch Abbildungen erläuterte Beschreibung einer Vorrichtung, mittelst welcher jenes Verhältniss sich werde festsetzen lassen, und zählt zuletzt die mancherley wichtigen Vortheile auf, welche die Astronomie und Physik von einer sichern und wichtigen Photometrie zu erwarten haben. --- Die in Vorschlag gebrachte Vorrichtung kommt im Wesentlichen mit dem seit vielen Jahren bereits bekannten \textsc{Köhler}schen Photometer ganz überein, dessen unmittelbare Anwendung nicht sehr brauchbar befunden worden ist, und unterscheidet sich von diesem nur durch einen künstlichern Mechanismus. Ob der Verfasser aber eine solche Vorrichtung wirklich ausgeführt, und ihre Leistungen bey den verschiedenen Lichtabstufungen der Sterne untersucht habe, geht aus der Abhandlung nicht hervor, indem darin überall keine Erfahrungen und Beobachtungen angeführt worden sind, aus welchen sich solche erkennen liessen. Da also diese Abhandlung die beiden Hauptbedingungen der Aufgabe unberücksichtigt gelassen hat, so konnte ihr der Preis nicht zuerkannt werden.

Der Verfasser der zweyten Concurrenzschrift hat dagegen beide Forderungen erfüllen gestrebt, indem er nicht nur das Lichtverhältniss auf eine neue, noch nicht versuchte, sinnreiche Art mittelst einer ganz einfachen, im Modell beygefügten, Vorrichtung zu bestimmen vorschlägt, sondern auch ihre Anwendbarkeit auf Sterne von ganz verschiedener Grösse durch das aus einer Reihe wirklich angestellter Beobachtungen erhaltene Resultat darzulegen sucht. Da aber theils diese Versuche zu wenig zahlreich und mannigfaltig sind, als dass dadurch schon eine Bürgschaft für die Richtigkeit der Resultate gewonnen wäre, auch eine eigentliche Theorie dieses Apparats, welche hier wesentlich und nothwendig ist, gänzlich vermisst wird: so konnte die Königl. Societät dieser Abhandlung in ihrer gegenwärtigen Gestalt den Preis nicht zuerkennen. In der Überzeugung jedoch, dass der Verfasser dieser Abhandlung, welcher Originalität, practischen Sinn und viele mechanische Anstelligkeit an den Tag legt, seinen Ideen mehr Vollendung gegeben, und die Anwendbarkeit des sinnreich gewählten Apparats durch zahlreichere Versuche mehr begründet haben würde, wenn die von ihm angeführten Hindernisse seine Musse zu diesen Untersuchungen weniger beschränkt hätten, findet die Kön. Societät
\par\HistoricalFoot{XI\textsubscript{1}.}{22}
\clearpage\SourceLeaf{175}\HistoricalHead{170}{L}{PHYSIK.\quad BRIEFWECHSEL.}
\noindent sich um so mehr bewogen, die Gelegenheit zur weitern Ausbildung derselben durch eine wiederholte Aufstellung dieser Preisfrage darzubieten, da sie die Hoffnung hegen darf, dass dadurch noch mehrere Mitbewerber veranlasst werden könnten, dieser der Wissenschaft so vielfachen Nutzen versprechenden Aufgabe ihre Untersuchungen zu widmen. Sie stellet daher diese Frage: »Über die Bestimmung der Lichtstärke der Himmelskörper« noch einmal für den November 1834 auf[*)].\begingroup\renewcommand{\thefootnote}{}\footnotetext{[*) Der Bericht über das Ergebnis dieser zweiten Preisausschreibung findet sich in den Göttingischen Gelehrten Anzeigen 1835, S. 329 ff., Werke VI, S. 649. Danach erhielt \textsc{Steinheil} den Preis, während das Accessit einer andern Arbeit zuerkannt wurde, als deren Verfasser sich (siehe Göttingische Gelehrte Anzeigen 1835, S. 1329) \textsc{Gerling} bekannt hat.]}\endgroup\ . . . . . . . .
\par\vspace{2\baselineskip}\begin{center}\rule{.24\linewidth}{.4pt}\end{center}
\clearpage\endgroup
